\section{Related Works}

\noindent\textbf{GraphRAG.} GraphRAG extends standard RAG by organizing corpus knowledge into a graph and performing retrieval over structured neighborhoods/paths/communities rather than isolated text chunks~\cite{zhang2025survey,nguyen2026urag}. Existing GraphRAG methods differ mainly in how they build and exploit this structure: tree-based indexing for coarse-to-fine retrieval (e.g., RAPTOR~\cite{sarthi2024raptor}); passage-graph variants that connect chunks via entity or similarity links (e.g., KGP~\cite{wang2024knowledge}); knowledge-graph approaches that extract entities/relations and reason over KGs or graph encoders (e.g., G-Retriever~\cite{he2024g}, HippoRAG~\cite{jimenez2024hipporag}, GFM-RAG~\cite{luo2025gfm}, KG$^2$RAG~\cite{zhu2025knowledge}, Graphusion~\cite{yang2025graphusion}, GNN-RAG~\cite{mavromatis2025gnn}); and rich-graph systems that augment graphs with attributes, summaries, or multi-granular indexing (e.g., Microsoft GraphRAG~\cite{edge2024local}, LightRAG~\cite{guo2024lightrag}, KET-RAG~\cite{huang2025ket}). 
While these systems target improved reasoning, moving from flat chunks to graph-structured retrieval also shifts the privacy attack surface: attackers may aim to reconstruct the underlying graph structure itself. 
More related work of GraphRAG is given in Appendix~\ref{graphrag-appendix}.

\noindent\textbf{Data Extraction in RAG and GraphRAG.} A growing body of work studies data extraction attacks against RAG systems. Prior work has shown that crafted prompts can induce verbatim leakage from retrieval stores~\cite{zeng2024good,qifollow,qi2026benchmarking}, while other attacks use poisoned knowledge bases~\cite{chaudhari2024phantom} or escalate from inference attacks to near-complete corpus reconstruction~\cite{cohen2024unleashing}. More recent approaches employ adaptive, agent-style extraction using anchor-based relevance search~\citealp{di2024pirates}, memory-focused probing~\citealp{wang2025unveiling,jiang2025feedback,wang2025silent}, and jailbreak-style prompt optimization~\cite{he2025external}.

\noindent As RAG systems increasingly incorporate structured intermediates, privacy analysis has also extended to GraphRAG. Graph-augmented LLM frameworks already suggest that gains come not only from retrieving isolated text units but also from explicitly reasoning over graph structure itself~\cite{jin2024graph}, which makes the security implications of graph-centered retrieval increasingly important. \citet{liu2025exposing} identifies a trade-off: GraphRAG may reduce verbatim leakage but can be more vulnerable to structured entity--relation extraction. More recent work further shows that closed-box, multi-turn subgraph reconstruction is feasible under realistic safeguards~\cite{song2026subgraph}, while GraphRAG-specific defenses have begun to emerge for protecting proprietary knowledge graphs against private-use theft~\cite{wang2026making}. However, existing work still provides a limited understanding of query-efficient, automatic recovery of the broader latent graph under a fixed attack budget, which is the setting we study here. A broader review is in Appendices~\ref{data_extraction_appendix} and~\ref{privacy_appendix}.

% \noindent\textbf{Data Extraction in RAG and GraphRAG.} A growing body of work studies data extraction attacks against RAG systems. Prior work has shown that crafted prompts can induce verbatim leakage from retrieval stores~\cite{zeng2024good,qifollow}, while other attacks use poisoned knowledge bases~\cite{chaudhari2024phantom} or escalate from inference attacks to near-complete corpus reconstruction~\cite{cohen2024unleashing}. More recent approaches employ adaptive, agent-style extraction using anchor-based relevance search~\citealp{di2024pirates}, memory-focused probing~\citealp{wang2025unveiling,jiang2025feedback,wang2025silent}, and jailbreak-style prompt optimization~\cite{he2025external}.

% As RAG systems increasingly incorporate structured intermediates, privacy analysis has also extended to GraphRAG. Graph-augmented LLM frameworks already suggest that gains come not only from retrieving isolated text units but also from explicitly reasoning over graph structure itself~\cite{jin2024graph}, which makes the security implications of graph-centered retrieval increasingly important.
% \citet{liu2025exposing} find a trade-off: GraphRAG may reduce verbatim leakage but can be more vulnerable to structured entity–relation extraction. However, existing attacks largely optimize leakage success rather than query-efficient reconstruction, and no study has extracted the latent graph under a fixed query budget—central to our setting. A broader review (privacy-aware and defensive RAG) is in Appendices~\ref{data_extraction_appendix} and~\ref{privacy_appendix}.

% We next review data extraction attacks on standard RAG.

% Long Version:
%Graph Retrieval-Augmented Generation (GraphRAG) extends standard RAG by building an explicit \emph{knowledge graph} over the corpus and using this structure to guide retrieval and reasoning~\cite{zhang2025survey}. Compared to vanilla RAG, which stores external knowledge as a flat collection of text chunks and relies on vector search, GraphRAG (i) represents knowledge as nodes and edges over entities, passages, and their relations, and (ii) performs retrieval over graph structures (e.g., neighborhoods, paths, or communities) rather than isolated chunks.
% Existing GraphRAG methods mainly differ in how they construct and exploit this graph. 
% Tree-based approaches such as RAPTOR~\cite{sarthi2024raptor} build hierarchical cluster trees for coarse-to-fine retrieval; 
% passage-graph methods like KGP~\cite{wang2024knowledge} link chunks into a passage graph via entity or similarity edges; 
% knowledge-graph methods such as G-Retriever~\cite{he2024g}, HippoRAG~\cite{jimenez2024hipporag}, GFM-RAG~\cite{luo2025gfm}, 
% KG$^2$RAG~\cite{zhu2025knowledge}, Graphusion~\cite{yang2025graphusion}, and GNN-RAG~\cite{mavromatis2025gnn} 
% explicitly extract entities and relations or leverage graph neural encoders over KGs; 
% and rich-graph variants like Microsoft GraphRAG~\cite{edge2024local}, LightRAG~\cite{guo2024lightrag}, and KET-RAG~\cite{huang2025ket} 
% augment graphs with textual attributes, summaries, or multi-granular indexing to support multi-scale retrieval.

% While these systems are primarily designed to improve reasoning quality, the move from flat chunks to structured graphs also changes the privacy attack surface: an adversary may now target the underlying graph itself rather than individual passages. We therefore next review data extraction attacks in standard RAG systems.

% \jiahao{Missing Citations (feel free not to add if you don't have time):
% \begin{itemize}
%     \item Knowledge graph-guided retrieval augmented generation (NAACL'25)~\cite{zhu2025knowledge}: A combination of KG-based graph RAG and vanilla RAG
%     \item KET-RAG: A cost-efficient multi-granular indexing framework for graph-rag (KDD'25)~\cite{huang2025ket}: An acceleration of microsoft's GraphRAG
%     \item Graphusion: A RAG Framework for Scientific Knowledge Graph Construction with a Global Perspective (WWW Workshop'25)~\cite{yang2025graphusion}: Put this in KG methods and no need to be too detailed.
%     \item GNN-RAG: Graph Neural Retrieval for Efficient Large Language Model Reasoning on Knowledge Graphs (ACL Findings'25)~\cite{mavromatis2025gnn}: Use GNNs for KG-based GraphRAG
% \end{itemize}
% }  \jiahao{Done}

% \subsection{Data Extraction in RAG and GraphRAG}
% \label{subsec:rag_extraction}


% Privacy analysis has recently extended to GraphRAG, where \cite{liu2025exposing} find a trade-off: GraphRAG may reduce verbatim leakage but can be more vulnerable to structured entity–relation extraction. However, existing attacks focus primarily on leakage success rather than query-efficient reconstruction, and none address extracting the entire latent graph under a fixed query budget—a key difference from our setting. A more detailed review, including privacy-aware and defensive RAG, is in Appendices~\ref{data_extraction_appendix} and \ref{privacy_appendix}.

% \noindent\textbf{Data Extraction in RAG and GraphRAG}. A growing body of work studies data extraction attacks against RAG systems. \citealp{zeng2024good} and \citealp{qifollow} show that carefully crafted prompts (including prompt-injection-style inputs) can induce verbatim leakage from flat retrieval stores in black-box settings; Phantom~\cite{chaudhari2024phantom} demonstrates backdoored leakage via knowledge-base poisoning; and \citealp{cohen2024unleashing} show escalation from membership/entity inference toward near-complete corpus reconstruction, analyzing the effects of design choices such as embedding type and context length. More recent work develops \emph{adaptive} and agent-style extraction, including anchor-based relevance search~\citealp{di2024pirates}, memory-focused probing (MEXTRA~\cite{wang2025unveiling}, CopyBreakRAG~\cite{jiang2025feedback}, IKEA~\cite{wang2025silent}), and a unified framework SECRET that combines jailbreak-style prompt optimization with adaptive triggering~\cite{he2025external}.
% Recent work has begun to extend privacy analysis to GraphRAG systems. \cite{liu2025exposing} designed tailored probes and report a trade-off: compared to standard RAG, GraphRAG may reduce verbatim text leakage but can be more vulnerable to extracting structured entity–relation information. 
% However, existing attacks largely focus on leakage success rather than query-efficient reconstruction, and do not consider the problem of extracting the entire latent graph under a fixed query budget, which clearly differs from our attack setting in this paper. A more detailed review along with Privacy-Aware and Defensive RAG is in Apendix~\ref{data_extraction_appendix} and \ref{privacy_appendix}.

% Long Version:
% A growing body of work studies data extraction attacks against RAG systems. Early methods focus on inducing verbatim leakage of private chunks from flat retrieval stores. \citealp{zeng2024good} and \citealp{qifollow} show that carefully crafted or injected prompts can cause RAG models to regurgitate proprietary text, even in black-box settings. 
% Phantom~\cite{chaudhari2024phantom} shows that knowledge-base poisoning can implant a backdoor in RAG, enabling verbatim passage leakage when a trigger appears in the user query. 
% \citealp{cohen2024unleashing} further demonstrate that such attacks can escalate from membership or entity-level inference to near-complete reconstruction of a QA chatbot’s underlying corpus, and systematically analyze how design choices such as embedding type and context length affect extraction.
% More recent work develops \emph{adaptive} extraction strategies and agent-style attacks. For example, \citealp{di2024pirates} %\jiahao{Fixed a typo for the method name here.} 
% use an anchor-based relevance search to balance exploration and exploitation over a hidden datastore. Memory-focused attacks such as MEXTRA~\cite{wang2025unveiling}, CopyBreakRAG~\cite{jiang2025feedback}, and IKEA~\cite{wang2025silent} show that RAG agents with working memory or history-aware prompts can be systematically probed to extract stored user data and retrieved chunks. \cite{he2025external} formalize external data extraction attacks and propose SECRET, a unified framework that combines jailbreak-style prompt optimization with adaptive triggering to efficiently elicit retrieval across many RAG instances.

% Recent work began to extend the study of privacy risks to GraphRAG systems. Specifically, \citealp{liu2025exposing} provide an early empirical analysis, designing tailored attacks to probe GraphRAG’s tendency to leak both raw text and structured information such as entities and relations. Their results reveal a trade-off: compared to standard RAG, GraphRAG may reduce verbatim text leakage but is more vulnerable to the extraction of structured entity-relationship information. However, existing attacks largely focus on leakage success rather than query-efficient reconstruction, and do not consider the problem of extracting the entire latent graph under a limited query budget, which clearly differs from our attack setting in this paper. 


% \noindent\textbf{Privacy-Aware and Defensive RAG} \suhang{Put this in appendix if we do not have enough space and mention it in main content}
% A complementary line of work designs privacy-preserving RAG pipelines. \citealp{zhou2025privacy} propose an encryption-based architecture that encrypts both text and embeddings before storage, restricting access to authorized clients with decryption keys. \citealp{mori2025differentially} introduce DP-SynRAG, which generates a differentially private \emph{synthetic} RAG database to improve utility under a fixed privacy budget. However, these defenses mainly target flat, text-chunk RAG and do not explicitly consider GraphRAG with structured graph-based knowledge bases, so their threat models and mechanisms may not transfer. We therefore study query-efficient graph extraction attacks on GraphRAG via a multi-agent, memory-augmented framework (AGEA).

%% Long Version %%
% A complementary line of work aims to design privacy-preserving RAG pipelines. \citealp{zhou2025privacy} propose an encryption-based architecture that encrypts both textual content and embeddings before storage, so only authorized clients with decryption keys can access the underlying data. 
% \citealp{mori2025differentially} introduce {DP-SynRAG}, which uses LLMs to generate a differentially private \emph{synthetic} RAG database that can be reused without repeated noise injection, improving utility under a fixed privacy budget. 
% %\jiahao{(Minor) Jiale's paper mainly focuses on showing the privacy risk of Graph RAG, instead of proposing a privacy-aware/defensive solution. It might be better to put it in the previous section Data Extraction in RAG.} \jiahao{Done}
% %For GraphRAG specifically, \cite{liu2025exposing} provides an initial empirical analysis of privacy risks. They design tailored attacks to probe GraphRAG’s tendency to leak both raw text and structured data such as entities and relations, and show a trade-off: compared to standard RAG, GraphRAG may reduce verbatim text leakage but is more vulnerable to extraction of structured entity–relationship information.

% These works focus on mitigating privacy leakage in standard RAG pipelines via cryptographic protection or differential privacy. However, they are primarily designed for text-chunk-based RAG and do not explicitly consider Graph RAG systems with structured, graph-based knowledge bases. As a result, their threat models and defense mechanisms may not directly apply to Graph RAG, motivating the need to study privacy risks and attacks tailored to graph-based RAG settings. 
% In this work, we focus on a multi-agent, memory-augmented framework, namely \emph{AGEA}, for query-efficient graph extraction attacks on GraphRAG systems, which are not covered by existing defense frameworks.

%These works seek to mitigate leakage in RAG and GraphRAG via cryptography, differential privacy, or high-level analyses of structured leakage. However, they do not model an adversary whose explicit goal is to reconstruct the \emph{entire latent knowledge graph} underlying a GraphRAG system, nor do they employ agentic querying over graph state and novelty. 
%In the next section, we address this gap by introducing \emph{AGEA}, a multi-agent, memory-augmented framework for query-efficient graph extraction attacks on GraphRAG systems.



% A complementary line of work focuses on designing privacy-preserving RAG pipelines. 
% Zhou et al.~\cite{zhou2025privacy} propose an encryption-based RAG architecture that encrypts both textual content and embeddings before storage. Only clients with the appropriate decryption keys can access the underlying data, reducing the risk of unintended data exposure from the retrieval datastore.

% Mori et al.~\cite{mori2025differentially} observe that existing private RAG methods typically rely on query-time differential privacy, which requires repeated noise injection and leads to accumulated privacy loss. They introduce \textbf{DP-SynRAG}, a framework that uses LLMs to generate a differentially private \emph{synthetic} RAG database. Once generated, the synthetic corpus can be reused without incurring additional privacy cost per query. DP-SynRAG extends private prediction by instructing LLMs to mimic subsampled database records in a DP manner, and experiments show that it outperforms prior private RAG methods under a fixed privacy budget.

% For GraphRAG specifically, Liu et al.~\cite{liu2025exposing} provide a first empirical analysis of privacy risks. They design tailored extraction attacks to probe GraphRAG’s tendency to leak both raw text and structured data such as entities and relations. Their study reveals a critical trade-off: compared to standard RAG, GraphRAG may reduce verbatim text leakage but becomes significantly more vulnerable to extraction of structured entity–relationship information, and they discuss several high-level defense directions.

% These works aim to \emph{mitigate} privacy risks in RAG and GraphRAG, either via cryptographic protection, differential privacy, or empirical analysis of leakage patterns. In contrast, our work focuses on a stronger attacker who seeks to \emph{reconstruct the entire latent knowledge graph} underlying a GraphRAG system. Our graph extraction attack thus serves as a complementary stress test for privacy-aware RAG designs, and can inform the development and evaluation of future defenses.

% Despite these defenses, existing work either protects flat text stores or provides high-level analyses of GraphRAG leakage, without modeling an adversary that explicitly targets the \emph{entire latent knowledge graph} or leverages agentic querying over graph state. In the next section, we introduce \emph{GraphMiner}, a multi-agent, memory-augmented framework for query-efficient graph extraction attacks on GraphRAG systems.


% In this work, we focus on GraphRAG systems that expose a graph-based retrieval interface but treat the underlying graph as a hidden internal state. Our goal is not to improve GraphRAG quality, but to study how an adversary can \emph{extract} this latent knowledge graph via black-box querying, across representative instantiations such as Micros